\documentclass[a4paper,10pt]{article}
\usepackage[left=3cm,top=3cm,right=2.5cm,bottom=4cm]{geometry}
\usepackage[utf8]{inputenc}
\usepackage[activeacute,spanish]{babel}
\usepackage{multirow}	% Para poder unir filas en las tablas
\usepackage{colortbl}	% Para colorear tablas
\usepackage{latexsym,amssymb,oldlfont,amsmath}
\usepackage{graphicx,graphics}
\usepackage{multicol}
\usepackage{color}
\usepackage{listings}
\usepackage{lastpage}
\usepackage{enumitem}



%\usepackage{fontspec}
%\setmainfont{Times New Roman}
\newcommand{\R}{{\mathbb R}}
\newcommand{\si}{{\mathcal S}^n}
\renewcommand{\theequation}{\thesection.\arabic{equation}} % Cambia formato salida numeracion equation
\makeatletter\@addtoreset {equation}{section} \makeatother % Tambien se usa
\renewcommand{\baselinestretch}{1.3}
\usepackage{amssymb}
\usepackage{hyperref}
\definecolor{midnightblue}{rgb}{0.1,0.1,0.44}


\newcommand{\dis}{\displaystyle}
\renewcommand{\theenumii}{\arabic{enumii}}

\def\Z{\mathbb{Z}}
\def\N{\mathbb{N}}
\def\R{\mathbb{R}}
\def\M{\mathcal{M}_n(\mathbb{R})}
\def\sin{\operatorname{sen}}
\def\e{\operatorname{e}}
\DeclareMathOperator{\arcsec}{arcsec}
\DeclareMathOperator{\arccot}{arccot}
\DeclareMathOperator{\arccsc}{arccsc}



\usepackage{etoolbox,fancyhdr,xcolor}
\fancyhf{} % sets both header and footer to nothing
\renewcommand{\headrulewidth}{0pt}
%\renewcommand{\footrulewidth}{0.5pt}
\pagestyle{fancy}
\headheight=60pt %para cambiar el tamaño del encabezadox
\fancyhead[L] %la "L" indica a la izquierda
{
\begin{minipage}{6cm}
\includegraphics[width=0.9\textwidth]{logo_udp.eps}
\end{minipage}
}

%\fancyhead[C]
%{
%\textbf{Cálculo II - Guía 1}
%}


%\fancyhead[R] %la "L" indica a la izquierda
%{

%}
%\begin{minipage}{4cm}
%\small %tamaño de letra pequeño
%{
%\textsf
%{
%Universidad de Chile\\
%...}}
%\end{minipage}}



%\fancyfoot[C]{}
%{
%\begin{minipage}{10cm}
%\color{midnightblue}
%\begin{center}Casa Matriz: Monjitas 527 Piso 15 - Santiago de Chile\\
%Fono (02)24367255 \end{center}
%\end{minipage}
%}
\fancyfoot[C]{}
%Página \thepage\ de \pageref{LastPage}}% Páginas Pares numeradas a la derecha
%\fancyfoot[LO]{\thepage} % Páginas Impares numeradas a la izquierda

\title{C2_Ay1_Ej2}

\begin{document}

\begin{center}
Cálculo II\\
Ayudantía Nº1 - Ejercicio Nº2\\
Primer Semestre 2017
\end{center}





2. Use integraci\'{o}n por partes para calcular las siguientes
integrales.
\begin{enumerate}[label=\emph{\alph*)}]
\item $\dis \int \cos(x) \ln(\sin(x))\, dx$
\item $\dis \int x^5\cos(x^3)\, dx$
\end{enumerate}


\vspace{5mm}

Solución:

\begin{enumerate}[label=\emph{\alph*)}]
\item $\dis \int \cos(x) \ln(\sin(x))\, dx$\\

Usando integración por partes:
\begin{align*}
 u  &=  \ln(\sin(x))  &   v  &= \sin(x) \\
 du &= \frac{\cos(x)}{\sin(x)} \, dx     &   dv &= \cos(x) \, dx
\end{align*}

Obtenemos lo siguiente:
\begin{align*}
 \dis \int \cos(x) \ln(\sin(x))\, dx &= \sin(x) \ln(\sin(x)) - \dis \int \frac{\cos(x)}{\sin(x)}\sin(x) \, dx \\
      &=  \sin(x) \ln(\sin(x)) - \dis \int \cos(x) \, dx \\
       &= \sin(x) \ln(\sin(x)) - \sin(x) +C
\end{align*}
\textcolor{red}{1 punto por desarrollar completa la a)}

\item $\dis \int x^5\cos(x^3)\, dx$\\
Usando integración por partes:
\begin{align*}
 u  &=  x^3  &   v  &= \frac{\sin(x^3)}{3} \\
 du &= 3x^2 \, dx     &   dv &= x^2\cos(x^3) \, dx
\end{align*}

Obtenemos lo siguiente:

\begin{align*}
 \dis \int x^5\cos(x^3)\, dx &= \frac{x^3\sin(x^3)}{3} - 3 \dis \int x^2 \frac{\sin(x^3)}{3} \, dx
 \end{align*}
 \textcolor{red}{1 punto por aplicar integración por partes}
 \begin{align*}
      &=  \frac{x^3\sin(x^3)}{3} - \dis \int x^2 \sin(x^3) \, dx \\
       &=  \frac{x^3\sin(x^3)}{3} - (- \frac{\cos(x^3)}{3}) + C \\
       &=  \frac{x^3\sin(x^3)}{3} + \frac{\cos(x^3)}{3} + C
\end{align*}
\textcolor{red}{1 punto por llegar al resultado final}
\end{enumerate}




\end{document}
